\documentclass[11pt]{article}
\usepackage{amsmath}
\usepackage{graphicx}
\usepackage{hyperref}

\title{Boot Claude: A Language Model Writes an Operating System and the Toolchain That Builds It}
\author{Martin Monperrus \and Claude}

\begin{document}
\maketitle

\begin{abstract}
We report on \texttt{boot-to-claude}, an x86 operating system written by Claude, a
large language model working as a coding agent. The system boots in QEMU through its own
two-stage boot loader into a 32-bit kernel with a disk driver, a file system, an
interactive shell, a scripting language, and a compiler to an in-kernel bytecode VM.
In a second stage, the agent replaced \texttt{nasm}, \texttt{gcc}, \texttt{ld} and
\texttt{objcopy} with its own assembler, C compiler and linker, and bootstrapped them
until the second and third self-compiled generations were byte-identical. The OS was
built in one evening; the toolchain spanned four calendar days. Commit messages document 43 defects
found and fixed along the way, almost all of them caught by booting the system rather
than by reading code.
\end{abstract}

\section{Introduction}\label{sec:intro}

Benchmarks for code-generating models have grown from single functions~\cite{chen2021codex}
to repository-level issue fixing~\cite{jimenez2024swebench}. Both measure changes to code
that already works. Writing a system from an empty repository is a different task. The
agent has to choose an architecture, keep it coherent over thousands of lines, and debug
failures whose only symptom may be a machine that prints nothing.

An operating system is a demanding case of this task. Unix showed that a small team
can build one that is both small and useful~\cite{ritchie1974unix}, and teaching kernels
such as xv6~\cite{cox2006xv6} show how much fits in a few thousand lines. Below the
kernel there is no runtime to report a mistake. A wrong descriptor table entry resets the
machine, and an I/O instruction with the wrong operand size corrupts the disk silently.

This paper reports what Claude built, in what order, and which defects it hit. The
source is public.\footnote{\url{https://github.com/monperrus/boot-to-claude}} All numbers
below are computed from its git history at commit \texttt{a4a64cb}.

\section{The System}\label{sec:system}

\paragraph{Boot and kernel.} The BIOS loads a 512-byte boot sector, which loads a second
stage. The second stage enables the A20 line, loads the kernel, installs a flat GDT and
switches to 32-bit protected mode. The kernel owns its GDT and an IDT that covers CPU
exceptions only. It has a free-list allocator over a fixed 1\,MB heap, VGA text and
COM1 serial output, a polled ATA PIO disk driver, and a polled PS/2 keyboard driver.

\paragraph{Deliberately single-tasking.} The kernel never enables interrupts after boot.
It has no scheduler, no paging and no user mode. The shell and both language runtimes run
in kernel context. This was a design decision stated in the repository: it keeps the OS
layer small and moves the interesting work into the file system and the languages.

\paragraph{Above the kernel.} \emph{tinyfs} is a flat file system with a superblock, a
free-block bitmap, a fixed inode table and 4\,KB maximum file size. Its format code is
compiled twice, once into the kernel and once into the host tool that formats the disk
image, so there is a single implementation of the on-disk format. The shell offers
\texttt{ls}, \texttt{cat}, \texttt{write}, \texttt{rm}, \texttt{run}, \texttt{build} and
\texttt{exec}. \texttt{run} executes a dynamically typed script with a tree-walking
interpreter. \texttt{build} compiles a statically scoped, integer-only C-like language to
stack bytecode, and \texttt{exec} runs it in a VM that supports recursion.

\paragraph{Toolchain.} \texttt{toolchain/} holds an assembler for the NASM subset the OS
uses, a compiler for a C dialect called tinyC, and a two-mode ELF32 linker. The host
\texttt{gcc} is invoked exactly once, to compile the first generation of this toolchain.
Table~\ref{tab:size} gives the size of each part.

\begin{table}[h]
\centering
\begin{tabular}{lr}
Component & Lines \\
\hline
Boot loader (\texttt{boot/}) & 218 \\
Kernel (\texttt{kernel/}) & 5003 \\
\quad of which languages (\texttt{kernel/lang/}) & 2041 \\
\quad of which game (\texttt{kernel/game/}) & 748 \\
\quad of which tinyfs (\texttt{kernel/fs/}) & 333 \\
\quad of which shell (\texttt{kernel/shell/}) & 286 \\
Toolchain (\texttt{toolchain/}) & 6430 \\
\quad of which C compiler & 3109 \\
\quad of which assembler & 2050 \\
\quad of which linker & 816 \\
Build and test tools (\texttt{tools/}) & 682 \\
Test fixtures (\texttt{tests/}) & 424 \\
Design documents (\texttt{docs/}) & 1099 \\
\end{tabular}
\caption{Size of \texttt{boot-to-claude} at commit \texttt{a4a64cb}, in physical lines
including comments.}
\label{tab:size}
\end{table}

\section{How It Was Built}\label{sec:history}

The history has 23 commits over six days, in three episodes (Table~\ref{tab:timeline}).
Each phase ends in a commit whose message states what was verified and how.

\begin{table}[h]
\centering
\begin{tabular}{llrr}
Commit time (2026, CEST) & Phase & Files & Lines added \\
\hline
Jul 17, 19:04 & 0: boot sector, build pipeline & 3 & 71 \\
Jul 17, 19:07 & 1: two-stage loader, protected mode & 10 & 350 \\
Jul 17, 19:10 & 2: GDT/IDT, allocator, self-test & 14 & 486 \\
Jul 17, 19:16 & 3: ATA driver, tinyfs & 13 & 620 \\
Jul 17, 19:20 & 4: keyboard driver, shell & 9 & 411 \\
Jul 17, 19:30 & 5: scripting language & 16 & 1126 \\
Jul 17, 19:38 & 6: compiler, bytecode VM & 13 & 1248 \\
Jul 17, 22:12 & 7: README, docs, final check & 2 & 181 \\
\hline
Jul 17, 23:40 & T1: assembler & 9 & 1823 \\
Jul 18, 08:36 & T2/T4: C compiler & 21 & 2999 \\
Jul 18, 08:49 & T3: adapt kernel to tinyC & 18 & 357 \\
Jul 18, 16:54 & T5: linker & 15 & 1101 \\
Jul 18, 17:11 & T6: drop \texttt{objcopy} & 4 & 66 \\
Jul 18, 22:59 & T7: switch the build over & 20 & 1022 \\
Jul 21, 07:17 & T8: self-hosting bootstrap & 30 & 1131 \\
Jul 21, 07:28 & T9: docs, cleanup & 1 & 23 \\
\hline
Jul 22, 07:20 & serial-playable battle game & 15 & 664 \\
Jul 22, 16:56 & VGA mode 13h driver, bitmap font & 5 & 428 \\
Jul 22, 17:21 & graphical battle screen & 10 & 520 \\
\end{tabular}
\caption{Phase commits of \texttt{boot-to-claude}. Four small commits (initial commit, a
GUI run target, a NASM macro expansion, a \texttt{.gitignore} entry) are omitted.}
\label{tab:timeline}
\end{table}

\paragraph{The OS in one evening.} Phases 0 to 6, from an empty boot sector to a
compiler with a bytecode VM, were committed within 34 minutes. Commit times are an upper
bound on the visible pace: they say when work was committed, not when it started.

\paragraph{The toolchain over four calendar days.} Replacing the host tools took longer than
writing the OS. The assembler had to grow from the instructions in four hand-written
\texttt{.S} files to everything the C compiler emits. The kernel source was first
rewritten into the tinyC dialect while still building with \texttt{gcc} (phase T3), so
that each change could be checked as behavior-preserving before the new toolchain was in
the loop.

\paragraph{The fixed point.} Phase T8 compiles the toolchain three times: Stage~1 with
\texttt{gcc}, Stage~2 with Stage~1, and Stage~3 with Stage~2. Let $S_i$ be the binaries of
stage $i$. The bootstrap is accepted only when
\begin{equation}\label{eq:fixpoint}
  S_2 = S_3 \quad \text{byte for byte, for each of } \texttt{as},\ \texttt{cc},\ \texttt{ld}.
\end{equation}
Equation~\eqref{eq:fixpoint} shows that the compiler compiles itself to the same result
it was compiled to. Getting there required features the toolchain had never needed for
the kernel: variadic functions, multi-dimensional arrays, and a minimal C library
written for the purpose.

\section{Defects and How They Were Found}\label{sec:defects}

Commit messages list each defect fixed in their phase. Counting them gives 43:
one in phase 3, two in phase 5, three in T1, three in T2/T4, six in T5, sixteen in T7,
eleven in T8, and one in the game. Some are missing features rather than wrong code,
such as \texttt{switch} statements not yet supported by the compiler. Three patterns
stand out.

\paragraph{Most defects were silent.} Few caused a crash. An \texttt{in}/\texttt{out}
instruction encoded without its operand-size prefix read two 16-bit words per access
instead of one, which surfaced as file system superblock corruption. A missing
\texttt{\textbackslash b} escape made the compiler lex \texttt{'\textbackslash b'} as
the letter \texttt{b}, so every typed \texttt{b} in a shell command acted as a
backspace. In the scripting interpreter, integer addition fell through to a default
branch returning 0, so \texttt{i = i + 1} looped forever.

\paragraph{Defects hid behind other defects.} A three-operand \texttt{imul} that dropped
its immediate went unnoticed because code that wrote and read arrays used the same wrong
index formula, so the errors cancelled. It surfaced only when linker-produced code read
an array written by the kernel.

\paragraph{The system test caught what unit tests missed.} The assembler's
\texttt{.dq} directive computed its high 32 bits with a shift by 32, which x86 masks to a
shift by 0. No toolchain test used \texttt{.dq}. The boot loader's disk address packet
does, and only the full boot suite caught it. The test oracle throughout is the serial
console: QEMU runs headless, a boot-time self-test prints \texttt{SELFTEST PASS}, and a
script types shell commands into QEMU's monitor and compares the serial output against
expected files.

\section{Threats to Validity}\label{sec:threats}

Kernel and compiler code is abundant in training data, so the agent may have reproduced
known designs rather than invented them. The defect count relies on commit messages
written by the agent, so it counts what the agent chose to report. This paper does not
yet measure how much a human steered the work, nor the cost in tokens or time spent per
phase. Finally, the system runs on QEMU only; real hardware is less forgiving.

\section{Related Work}\label{sec:related}

HumanEval~\cite{chen2021codex} measures function-level synthesis, and
SWE-bench~\cite{jimenez2024swebench} measures issue resolution in existing repositories.
\texttt{boot-to-claude} is a single construction of a whole system from scratch.
xv6~\cite{cox2006xv6} is the natural size reference for a teaching kernel. Unlike
\texttt{boot-to-claude}, it is multi-processor and has processes and paging. seL4~\cite{klein2009sel4}
shows that a kernel can be built against a machine-checked specification. Its checks are
proofs. The checks here are tests that boot the machine. Thompson's Turing lecture on
trusting compilers~\cite{thompson1984trusting} is the classic warning that a self-hosted
toolchain is only as trustworthy as its first generation. Here that generation is built
by \texttt{gcc} from source the agent wrote.

\section{Conclusion}\label{sec:conclusion}

A language model wrote a small, single-tasking operating system, then wrote the
assembler, compiler and linker that build it, up to a byte-identical self-hosting fixed
point. The OS was the easy part. Most of the effort and most of the defects went into the
toolchain, and most defects were found by booting the machine.

\bibliographystyle{plain}
\bibliography{refs}

\end{document}
